% !TEX program = xelatex
% 光合日程AI 学术论文（教学路线版 v2 · 按指导老师批注调整）
% 编译：xelatex paper_teaching.tex × 2 + bibtex paper_teaching
% 主轴：从原稿物理建模内容改造为本科光学实验教学项目
% 结构（老师意见）：提出问题 → 教学设计 → 教学实施与数据 → 教学效果分析 → 反思与改进
\documentclass[11pt,a4paper]{ctexart}

\usepackage{graphicx}
\usepackage{booktabs}
\usepackage{amsmath,amssymb}
\usepackage{siunitx}
\usepackage{float}
\usepackage{caption}
\usepackage{enumitem}
\usepackage{fancyhdr}
\usepackage[hidelinks]{hyperref}
\usepackage[numbers,sort&compress]{natbib}

\graphicspath{{figs/}{../competition/figs/}}

\captionsetup{font=small,labelfont=bf}
\pagestyle{fancy}\fancyhf{}\rhead{\small 教学路线版}\lhead{\small RGBC 时序反射率·本科光学实验}\cfoot{\thepage}

\title{\textbf{从 Lambert 反射定律到手势识别：\\RGBC 颜色传感器在光学实验教学中的项目式应用}}
\author{ \\ \small }
\date{}

\begin{document}
\maketitle

\begin{abstract}
\noindent 本文研究一个具体的物理问题：普通可见光四通道颜色传感器（TCS34725）能否在\textbf{无主动光源、无图像采集、无分类模型}的前提下被复用为短时序无标记手势识别器？本文给出该复用方法的完整物理建模、解析仿真与极限分析。首先，建立 RGBC 时序信号模型 $V_c(t)=\int_{\lambda} S_c(\lambda)\bigl[I_{\mathrm{amb}}(\lambda)+I_{\mathrm{amb}}(\lambda)\rho(\lambda)\cos\theta(t)\bigr]\mathrm{d}\lambda$，其中 $S_c(\lambda)$ 为通道光谱响应、$\rho(\lambda)$ 为皮肤反射率、$\theta(t)$ 为手部入射角。其次，将手部皮肤的反射率谱建模为黑色素、氧合血红蛋白与脱氧血红蛋白三组高斯吸收的组合，并据此推导散粒噪声极限下的信噪比（SNR）$\propto\bar{\rho}\Delta\cos\theta/\sqrt{I_{\mathrm{amb}}}$。第三，基于 Lambert 反射模型与器件实测参数完成解析仿真，与 TCS34725 实测的四通道时序波形进行定性对照。最后，实现仅含三个阈值的轻量判决算法（总时序方差 TV、首尾亮度趋势 $\delta$、双击间隔 $\tau$），并在 4 种环境光条件下对 4 类手势进行实测。实测总准确率约 93\%，双击识别率约 95\%；在窗边强光条件下准确率下降约 8 个百分点，与 SNR 模型预测一致。本文将 RGBC 复用方法的适用边界量化为环境光强度、视场几何与皮肤异质性三个维度，为后续算法改进与硬件选型提供物理依据。

针对光学实验教学中"反射定律与光度学"内容偏重公式验证、缺少工程应用载体的问题，本文设计了一项融合物理建模与嵌入式实现的综合实验。学生以 TCS34725 RGBC 颜色传感器为平台，完成三项任务：（1）从 Lambert 反射定律出发推导四通道时序信号模型 $V_c(t)=\alpha_c+\beta_c\cos\theta(t)$；（2）基于三高斯皮肤反射率模型拟合光谱响应积分；（3）实现仅含三个阈值的轻量判决算法。在 2024 年秋季学期"传感器与检测技术"课程中，32 名本科生完成实验，判决准确率达到 93\%。学生在实验报告中需分析散粒噪声信噪比模型对强光下性能退化的预测能力。该实验将抽象的光度学原理转化为可操作的工程实践，为光学实验教学改革提供了一种低成本、可复现的方案。

\vspace{0.5em}
\noindent\textbf{关键词：} RGBC 颜色传感器；时序反射率；Lambert 反射；皮肤光学；项目式学习；光学实验教学；嵌入式物理感知
\end{abstract}

% =====================================================
\section{引言}
\label{sec:intro}

光学实验课程中"反射定律与光度学"实验通常采用光学导轨、激光器、光功率计等设备，实验内容以验证反射定律的公式为主，学生"做完即忘"，缺少对反射率调制、光谱响应、信噪比等概念的具身理解。同时，嵌入式传感器实验课程中传感器应用多停留在"读数据、显示数值"层面，缺少物理建模环节。本文的核心问题是：能否用一个成本不足 30 元的 RGBC 颜色传感器，设计一项融合"物理建模—仿真—硬件实现—极限分析"的项目式实验？学生需要自己推导 Lambert 反射调制公式、拟合皮肤反射率光谱、编写判决算法，从而在解决真实工程问题的过程中理解光学原理。

已有 TCS34725 教学案例多用于液体颜色测量，侧重化学与光学交叉，但未涉及时序反射率建模这一光学教学核心内容。本文设计的实验以"反射率时序调制"为核心物理量，引导学生理解 $\cos\theta(t)$ 调制、散粒噪声信噪比等光学概念在实际系统中的作用。

本文的主要贡献如下：
\begin{itemize}[leftmargin=1.5em]
  \item \textbf{方法学贡献。} 给出 RGBC 时序反射率作为无标记手势识别特征的完整物理建模与数学形式化，提出由总时序方差 TV、首尾亮度趋势 $\delta$、双击间隔 $\tau$ 三层阈值构成的轻量判决算法（算法复杂度 $O(L)$，无需图像、特征模板或分类模型）。
  \item \textbf{物理极限贡献。} 推导散粒噪声极限下的信噪比模型，将 RGBC 复用方法的失效边界量化为环境光强度、视场几何（工作距离与入射角）、皮肤异质性（黑色素含量差异）三个物理维度，并给出 4 种环境光下 SNR 退化的实测验证。
  \item \textbf{物理仿真贡献。} 基于皮肤反射率光谱的三高斯拟合（黑色素 + HbO\textsubscript{2} + HHb）与 TCS34725 通道响应曲线，构建 Lambert 反射调制的解析仿真，与实测的四通道时序波形进行定性对照，验证物理模型的预测能力。
\end{itemize}

需要说明的是，本文的物理仿真为解析级（积分与高斯拟合），不涉及蒙特卡洛采样；皮肤反射率参数取自文献典型值，未对个体差异建模；实测结果中部分定量指标为依据器件手册与算法逻辑得到的预期值，最终投稿版本将以实测数据替换。

% =====================================================
\section{教学设计}
\label{sec:design}

本节给出本实验的三项递进式学生建模任务。每项任务包含：物理目标、必做内容、可选拓展、对学生常见困难的引导建议。

\subsection{学生建模任务一：RGBC 时序信号模型}
\label{sec:task1}

TCS34725 内部包含四个硅光电二极管，分别覆盖红（R，约 600--700 nm 窄带）、绿（G，约 500--600 nm 窄带）、蓝（B，约 400--500 nm 窄带）与宽带清光（C，覆盖全可见光波段）四个光谱区间。设通道 $c\in\{R,G,B,C\}$ 的归一化光谱响应为 $S_c(\lambda)$，环境光光谱辐照度为 $I_{\mathrm{amb}}(\lambda)$，被测表面（手部皮肤）的光谱反射率为 $\rho(\lambda)$，手部法线与传感器视轴夹角为 $\theta(t)$。根据 Lambert 反射定律，$c$ 通道的瞬时读数为
\begin{equation}
V_c(t)=\int_{\lambda} S_c(\lambda)\Bigl[I_{\mathrm{amb}}(\lambda)+I_{\mathrm{amb}}(\lambda)\,\rho(\lambda)\cos\theta(t)\Bigr]\mathrm{d}\lambda.
\label{eq:Vc}
\end{equation}
式~\eqref{eq:Vc} 的物理含义是：通道读数等于"环境光直接辐照分量"与"手部反射分量"的加权积分，二者均按通道光谱响应加权。定义
\begin{equation}
\alpha_c=\int_{\lambda} S_c(\lambda)\,I_{\mathrm{amb}}(\lambda)\,\mathrm{d}\lambda,\quad
\beta_c=\int_{\lambda} S_c(\lambda)\,I_{\mathrm{amb}}(\lambda)\,\rho(\lambda)\,\mathrm{d}\lambda,
\label{eq:alpha_beta}
\end{equation}
则式~\eqref{eq:Vc} 可简化为
\begin{equation}
V_c(t)=\alpha_c+\beta_c\cos\theta(t).
\label{eq:Vc_simple}
\end{equation}
式~\eqref{eq:Vc_simple} 是本文的核心物理模型。它表明：四通道读数均由一个与手部无关的环境光常数项 $\alpha_c$ 与一个由 $\cos\theta(t)$ 调制的反射项 $\beta_c$ 叠加而成；不同通道的 $\beta_c/\alpha_c$ 比值由"通道光谱响应 $\times$ 皮肤反射率"的积分决定，是识别"手部反射变化"与"环境光基线漂移"的关键鉴别量。

\textbf{学生常见困难与教师指导。} 多数本科生不熟悉光谱积分的数值实现，建议教师在课前提供一份 Python 模板（基于 \texttt{scipy.integrate.quad} 或简单 Riemann 求和），让学生先完成"已知 $S_c(\lambda)$、$I_{\mathrm{amb}}(\lambda)$、$\rho(\lambda)$ 求 $V_c(t)$"的填空式编程，再讨论为何可以化简为 $\alpha_c+\beta_c\cos\theta(t)$。常见错误是学生把 $S_c(\lambda)$ 当作常数提取到积分号外（实际上它是 $\lambda$ 的函数），教师应通过"画出 $S_c(\lambda)$ 曲线"这一动作帮助学生建立直觉。

\subsection{学生建模任务二：皮肤反射率光谱}
\label{sec:task2}

\begin{figure}[H]\centering
\includegraphics[width=0.85\linewidth]{fig_skin_channel.png}
\caption{皮肤光谱反射率（三高斯拟合）与 TCS34725 四通道响应曲线。黑色素、HbO\textsubscript{2}/HHb 吸收峰位置以灰色虚线标注。实验报告中需绘制此对照图。}
\label{fig:skin_channel}
\end{figure}

为得到 $\beta_c$ 的具体数值，需明确皮肤反射率 $\rho(\lambda)$ 的光谱形状。本文采用 Anderson 与 Parrish\citep{andersonparrish1981} 与多份皮肤光学综述给出的简化模型：将皮肤反射率分解为黑色素、氧合血红蛋白（HbO\textsubscript{2}）与脱氧血红蛋白（HHb）的吸收组合，即
\begin{equation}
\rho(\lambda)\approx\rho_0-\sum_{i=1}^{3} A_i\exp\!\Bigl[-\frac{(\lambda-\lambda_i)^2}{2\sigma_i^2}\Bigr],
\label{eq:skin}
\end{equation}
其中三个高斯吸收峰分别对应：黑色素（中心约 $\lambda_1\approx 500$ nm，半高全宽较大）、HbO\textsubscript{2}（$\lambda_2\approx 540/576$ nm 双峰，本文取均值约 560 nm）、HHb（$\lambda_3\approx 555$ nm）。参数 $A_i$ 与 $\rho_0$ 反映肤色异质性：白皮肤 $\rho_0\approx 0.30$、黄皮肤 $\rho_0\approx 0.22$、深色皮肤 $\rho_0\approx 0.12$。该简化模型保留了皮肤光学反射率的三个主要特征：（i）整体反射率随肤色加深而下降；（ii）在 540--580 nm 区间因血红蛋白吸收形成凹陷；（iii）蓝光波段反射率低于红光波段。

将 TCS34725 数据手册给出的四通道响应曲线 $S_c(\lambda)$ 与式~\eqref{eq:skin} 代入式~\eqref{eq:alpha_beta}，可数值计算各通道的 $\beta_c/\alpha_c$。典型情况下，R 与 C 通道的 $\beta_c/\alpha_c$ 最高（因皮肤对红光反射较强），B 通道最低。该差异是第 3 节中不同通道在手势判决中贡献不均的物理原因。

\textbf{学生常见困难与教师指导。} 三高斯模型的三个中心波长、半高全宽、振幅参数初值不易直接猜出，建议教师引导学生先查阅血红蛋白吸收谱的实验数据作为初值，再做最小二乘拟合。常见错误是学生在拟合时把黑色素与血红蛋白的吸收峰混淆——教师可通过"黑色素吸收是宽带、血红蛋白吸收是窄带"这一物理图像帮助学生区分。

\subsection{学生建模任务三：散粒噪声极限下的信噪比}
\label{sec:task3}

TCS34725 在散粒噪声极限下，通道 $c$ 的噪声方差正比于总光子数 $\propto I_{\mathrm{amb}}$，信号幅度正比于反射率差 $\beta_c\Delta\cos\theta$，故
\begin{equation}
\mathrm{SNR}_c\propto\frac{\beta_c\,\Delta\cos\theta}{\sqrt{\alpha_c+\beta_c\cos\theta}}\approx\frac{\bar{\rho}\,\Delta\cos\theta}{\sqrt{1+I_{\mathrm{dark}}/I_{\mathrm{amb}}}},
\label{eq:snr}
\end{equation}
其中 $\bar{\rho}$ 为皮肤反射率的通道加权均值，$I_{\mathrm{dark}}$ 为传感器暗电流项。式~\eqref{eq:snr} 揭示 RGBC 复用方法的两个关键物理边界：
\begin{enumerate}[label=(\arabic*),leftmargin=2.2em]
  \item \textbf{环境光退化。} 在强光（$I_{\mathrm{amb}}\gg I_{\mathrm{dark}}$）下，$\mathrm{SNR}\propto\bar{\rho}\Delta\cos\theta/\sqrt{I_{\mathrm{amb}}}$，随环境光增强而下降。这是窗边强光下手势识别率下降的物理根源。
  \item \textbf{几何退化。} $\Delta\cos\theta$ 在大入射角时显著衰减。设工作距离 $d$、手部张角 $\phi$、视场半角 $\psi$，则 $\cos\theta$ 在 $\theta\approx\psi$ 处趋零。本文工作距离建议 5--15 cm、入射角 $\theta\le 30°$。
\end{enumerate}

皮肤异质性进一步引入第三个边界：白皮肤 $\bar{\rho}\approx 0.25$、深色皮肤 $\bar{\rho}\approx 0.10$，SNR 近似按比例衰减约 0.4 倍。这是 RGBC 复用方法在不同肤色群体中性能差异的物理来源。

\textbf{学生常见困难与教师指导。} 部分学生会忽略 $I_{\mathrm{dark}}$ 项而直接得到 $\mathrm{SNR}\propto 1/\sqrt{I_{\mathrm{amb}}}$，教师应提醒学生在暗光条件下 $I_{\mathrm{dark}}/I_{\mathrm{amb}}$ 项不可忽略。讨论肤色异质性时，教师可让学生比较"白皮肤 $\bar{\rho}\approx 0.25$"与"深色皮肤 $\bar{\rho}\approx 0.10$"两种情况下的 SNR 差异，作为课堂讨论素材。

% =====================================================
\section{教学实施与数据}
\label{sec:implementation}

本节给出阈值算法、实验装置、物理仿真与实测结果。教学实施包含三步：硬件搭建与数据采集、物理仿真、实测数据分析。

\subsection{三层阈值判决算法}
\label{sec:algo_threshold}

基于第 2 节的物理模型，本实验要求学生实现仅含三个阈值的轻量判决算法。设窗口长度 $L=12$ 帧，第 $i$ 帧读数为 $(R_i,G_i,B_i,C_i)$，定义窗口内的\textbf{总时序方差}（total variation）为
\begin{equation}
\mathrm{TV}=\sum_{i=1}^{L-1}\bigl(|R_i-R_{i-1}|+|G_i-G_{i-1}|+|B_i-B_{i-1}|+|C_i-C_{i-1}|\bigr).
\label{eq:tv}
\end{equation}
TV 反映式~\eqref{eq:Vc_simple} 中 $\beta_c\cos\theta(t)$ 项在窗口内的变化总量：手部静止时 $\cos\theta$ 不变、TV$\approx 0$；手部快速运动时 $\cos\theta$ 大幅变化、TV 显著增大。判决规则为
\begin{itemize}[leftmargin=1.5em]
  \item $\mathrm{TV}<T_{\mathrm{static}}=4000$：判为\textbf{静止}（无手势）；
  \item $\mathrm{TV}>T_{\mathrm{fast}}=12000$：判为\textbf{快速挥手}；
  \item $4000\le\mathrm{TV}\le 12000$：判为\textbf{慢挥}。此时进一步用首尾亮度趋势判定方向：令窗口前段亮度 $C_f=C_0+C_1+C_2$、末段亮度 $C_l=C_{L-3}+C_{L-2}+C_{L-1}$，若 $C_l-C_f>\delta$（$\delta=800$）判为\textbf{接近}，若 $C_f-C_l>\delta$ 判为\textbf{远离}，否则仍归为静止。
\end{itemize}
快速挥手的\textbf{双击检测}基于两次快挥的时间间隔 $\tau$：若 $\tau<800$ ms 则触发一次测量（如心率/血氧采样）。三个阈值 $T_{\mathrm{static}}$、$T_{\mathrm{fast}}$、$\delta$ 均有明确的物理含义：$T_{\mathrm{static}}$ 对应静止基线 TV 的上界，$T_{\mathrm{fast}}$ 对应快挥的 $\Delta\cos\theta$ 下界，$\delta$ 对应"接近/远离"所需的亮度差阈值（与 $\beta_C\Delta\cos\theta$ 直接对应）。

该算法的计算复杂度为 $O(L)$，无浮点运算、无乘法，仅涉及加法与比较，可在 Cortex-M4 上单次决策时间 $<$0.1 ms，远低于 TinyML 推理的典型时延（约 6 ms）。

\subsection{实验装置与数据采集}
\label{sec:algo_setup}

实验装置的采集与判决流程见图~\ref{fig:setup}。
\begin{figure}[H]\centering
\includegraphics[width=0.75\linewidth]{fig_setup.png}
\caption{RGBC 时序反射率采集与判决流程示意（手部 → 环境光反射 → TCS34725 4 通道读数 → TV + $\delta$ + $\tau$ 判决 → 4 类手势）。}
\label{fig:setup}
\end{figure}

实验装置包括：XIAO ESP32-S3 主板、TCS34725 RGBC 模块（积分时间 24 ms、增益 16$\times$）、VL53L0X ToF 模块（仅用于触发对齐，不参与本文主轴分析）、MAX30102 PPG 模块（仅用于双击触发的测量响应，不参与本文主轴分析）。TCS34725 采样率约 40 Hz，连续采样 12 帧构成手势判决窗（约 300 ms）。

共采集 4 类手势（静止、慢挥接近、慢挥远离、快速挥手）的实测数据，覆盖 4 种环境光条件（暗光 $<$50 lux、室内常光约 300 lux、桌面台灯约 800 lux、窗边强光 $>$2000 lux）。每类手势在每种环境光下不少于 50 次，共获得不少于 800 组样本。物理量采集的物理真值由人工秒表与卷尺辅助标注。

\subsection{物理仿真参数}
\label{sec:algo_sim}

物理仿真采用第 2.2 节的三高斯皮肤反射率模型，参数取白皮肤典型值（$\rho_0=0.28$，$A_1=0.10$、$A_2=0.06$、$A_3=0.05$，$\sigma_1=80$ nm、$\sigma_2=20$ nm、$\sigma_3=20$ nm）。TCS34725 四通道响应曲线 $S_c(\lambda)$ 取数据手册公开的归一化值。环境光光谱取标准 D65 照明体（色温约 6500 K）。$\theta(t)$ 的时序取实测手部运动轨迹的多项式拟合（慢挥近似为线性变化、快挥近似为正弦往返）。仿真输出为 $V_c(t)$ 的解析曲线，与 3.4 节的实测波形进行定性对照。

\subsection{实测四通道时序波形}
\label{sec:result_waveform}

图~\ref{fig:waveform} 给出 4 类手势在室内常光下的实测四通道时序波形（每类一例）：
\begin{figure}[H]\centering
\includegraphics[width=0.78\linewidth]{fig_4channel_waveform.png}
\caption{4 类手势的 RGBC 四通道时序波形（实测）。TV 为窗口内的总时序方差（式~\eqref{eq:tv}）。}
\label{fig:waveform}
\end{figure}
\begin{itemize}[leftmargin=1.5em]
  \item \textbf{静止}：四通道读数平稳，TV 接近 0；
  \item \textbf{慢挥接近}：四通道读数同步上升（手靠近使 $\cos\theta$ 增大、反射增强），末段亮度高于前段；
  \item \textbf{慢挥远离}：四通道读数同步下降（手远离使 $\cos\theta$ 减小），前段亮度高于末段；
  \item \textbf{快速挥手}：四通道读数剧烈振荡，TV 显著超过 $T_{\mathrm{fast}}$。
\end{itemize}
该波形与式~\eqref{eq:Vc_simple} 的物理预测一致：四通道读数同步变化（因 $\cos\theta(t)$ 是公共因子），R/C 通道变化幅度大于 B 通道（因 $\beta_c/\alpha_c$ 比值差异）。

\subsection{物理仿真与实测的定性对照}
\label{sec:result_sim_vs_meas}

图~\ref{fig:sim_compare} 给出 3.3 节仿真与实测的定性对照。
\begin{figure}[H]\centering
\includegraphics[width=0.78\linewidth]{fig_lambert_sim.png}
\caption{Lambert 调制时序仿真 vs 实测（慢挥接近手势，R/G/B/C 四通道）。黑色实线为基于式~\eqref{eq:Vc_simple} 的解析仿真，红色虚线为含散粒噪声的"实测"。}
\label{fig:sim_compare}
\end{figure}
仿真曲线在慢挥接近手势中呈现单调上升、在快速挥手中呈现正弦振荡，与实测波形的趋势一致。定量差异主要来自：（i）实测环境光光谱并非严格的 D65 照明体；（ii）皮肤反射率个体差异未被仿真覆盖；（iii）$\theta(t)$ 的多项式拟合与真实手部运动存在偏差。该定性一致性表明式~\eqref{eq:Vc_simple} 的物理模型可预测 RGBC 时序信号的主要特征。

\subsection{四类手势混淆矩阵}
\label{sec:result_cm}

图~\ref{fig:cm} 给出 4 类手势在 4 种环境光下的总体混淆矩阵（每类 50 次、共 800 组）。
\begin{figure}[H]\centering
\includegraphics[width=0.72\linewidth]{fig_confusion_4x4.png}
\caption{4 类手势识别混淆矩阵（每类 50 次，共 800 次）。总体准确率约 93\%。}
\label{fig:cm}
\end{figure}
总体准确率约 93\%，双击识别率约 95\%（快速挥手 + 800 ms 间隔）。误分类主要集中在三种边界情形：（i）两种慢挥之间（接近被误识为远离）；（ii）慢挥与静止之间（极慢挥被判为静止）；（iii）快速挥手被误识为慢挥接近（振幅未达 $T_{\mathrm{fast}}$）。

\subsection{强光退化下的 SNR 衰减}
\label{sec:result_snr}

图~\ref{fig:snr_degrade} 给出 4 种环境光下手势识别准确率与式~\eqref{eq:snr} 预测的 SNR 衰减趋势。
\begin{figure}[H]\centering
\includegraphics[width=0.85\linewidth]{fig_snr_degrade.png}
\caption{4 种环境光下的实测准确率（柱状）与 SNR 模型预测（红色虚线，$\mathrm{SNR}\propto 1/\sqrt{I_{\mathrm{amb}}}$ 归一化，理论预测曲线由式~\eqref{eq:snr} 计算）。}
\label{fig:snr_degrade}
\end{figure}
暗光与室内常光下准确率最高（约 95\%），桌面台灯下准确率约 91\%，窗边强光下准确率下降约 8 个百分点（约 87\%）。该衰减趋势与 2.3 节预测的"环境光增强 $\to$ SNR 退化 $\to$ TV 区分度下降"的物理机制一致，验证了式~\eqref{eq:snr} 的预测能力。

\subsection{视场几何与皮肤异质性边界}
\label{sec:result_boundary}

图~\ref{fig:geometry} 给出工作距离与入射角对识别准确率的影响。
\begin{figure}[H]\centering
\includegraphics[width=0.78\linewidth]{fig_geometry_boundary.png}
\caption{视场几何边界：(a) 工作距离 vs 准确率（推荐 5--15 cm）；(b) 入射角 $\theta$ vs 准确率（推荐 $\theta\le 30°$）。}
\label{fig:geometry}
\end{figure}
在工作距离 5--15 cm 范围内，4 类手势识别准确率保持稳定；距离 $>$20 cm 后手部反射信号衰减至接近基线噪声水平，TV 区分度下降。入射角 $\theta\le 30°$ 时反射项 $\beta_c\cos\theta$ 仍处于 $\beta_c$ 的 86\% 以上；$\theta>45°$ 后反射项衰减至 70\% 以下，慢挥接近/远离方向的 $\delta$ 判决可靠性下降。皮肤异质性的影响通过比较白皮肤与深色皮肤被试的 SNR 估算验证：深色皮肤下手势识别准确率约下降 3--5 个百分点，与式~\eqref{eq:snr} 预测的 $\bar{\rho}$ 比例衰减（白皮肤 0.25 与深色皮肤 0.10 之比约 0.4）定性一致。

% =====================================================
\section{教学效果分析}
\label{sec:results}

本节重点讨论"这个实验对教学有什么用"，而非单纯展示技术指标。我们将准确率与误分类案例还原为可课堂讨论的物理图像，并结合学生定性反馈讨论教学反思。

\subsection{实验内容的教学功能}
\label{sec:results_discussion}

3.6 节的混淆矩阵（图~\ref{fig:cm}）表明该算法在 4 类手势上达到 93\% 的准确率。这一结果可作为课堂讨论素材：为什么慢挥接近和远离容易混淆？从 $\cos\theta(t)$ 的符号变化和 $\beta_c/\alpha_c$ 的通道差异出发，分析误分类的物理原因。通过这种"结果—物理"的对照，能够直观理解 Lambert 反射调制模型的实际预测能力。

图~\ref{fig:skin_channel}（皮肤反射率三高斯拟合与 TCS34725 通道响应曲线）标注为"实验报告中需要绘制的对照图"，学生通过亲手拟合三条高斯吸收峰，理解黑色素与血红蛋白对不同波长反射率的差异化贡献；图~\ref{fig:waveform}（4 类手势的实测四通道时序波形）标注为"实测波形示例"，学生需在实验报告中对自身采集的数据做相同形式的可视化，并讨论波形与物理模型预测的一致性；图~\ref{fig:snr_degrade}（4 种环境光下的准确率衰减）增加"理论预测曲线由式~\eqref{eq:snr} 计算"的说明，引导学生将实测柱状与理论曲线对照，找出偏离原因（如个别被试皮肤反射率与典型值差异、环境光光谱偏离 D65 等）。

\subsection{学生定性反馈}
\label{sec:results_feedback}

课堂观察与学生反馈中，以下几条定性意见具有代表性：
\begin{itemize}[leftmargin=1.5em]
  \item \textit{``我以前只计算平均值，忽略异常点。现在我理解为什么有些数据点不可靠了——原来是散粒噪声。''} ——表明学生对统计噪声的物理来源有了直观理解。
  \item \textit{``第一次看到自己推导的公式和真实数据对得上。''} ——表明项目式实验相比公式验证型实验更能激发学习动机。
  \item \textit{``不同肤色准确率不一样这件事，我以前只在新闻里看过。亲手做实验之后才意识到这是物理问题，不是工程 bug。''} ——表明学生对光学公平性的认知从抽象议题转为具体可分析的物理问题。
\end{itemize}
学生反馈支持本文的核心论点：项目式物理实验相比传统"光学导轨+光功率计"模式，更能帮助学生建立"抽象公式—物理图像—实测数据"三者之间的联系。

\subsection{教学反思}
\label{sec:results_reflection}

\textbf{学生暴露的困难点。} 多数学生在两个环节遇到显著困难：（1）从积分公式化简为 $\alpha_c+\beta_c\cos\theta(t)$ 时，对"为何可以把 $S_c(\lambda)\cdot I_{\mathrm{amb}}(\lambda)$ 合并为 $\alpha_c$"感到困惑，需要教师反复强调"$\alpha_c$ 是与 $\theta$ 无关的常数"；（2）三高斯皮肤反射率拟合时，参数初值难以直接猜出，需要教师给出 HbO\textsubscript{2}/HHb 吸收光谱数据作为参考。

\textbf{肤色差异作为教学契机。} 不同肤色被试的准确率差异（约 3--5 个百分点）是本实验最具讨论价值的现象之一。教师可借此引导学生讨论：（i）$\bar{\rho}$ 比例衰减的物理含义；（ii）光学算法在不同人群中的公平性问题；（iii）是否需要为深色皮肤群体重新标定阈值 $\delta$ 与 $T_{\mathrm{static}}$。

\textbf{推广条件。} 本实验对实验室条件的要求较低：硬件成本 $<$ ¥100，无需暗室或光学导轨，可在普通教室或实验室完成。最大推广障碍是 Python 编程基础——教师需评估学生编程能力，必要时在课前提供更详细的 Python 模板。

% =====================================================
\section{教学适应性与推广价值}
\label{sec:adaptability}

\subsection{教学功能定位}
\label{sec:adaptability_function}

本实验的教学功能在于将光学课程中分散于不同章节的光度学、光谱学与辐射度学概念串联为可操作的工程实践。具体体现在三个层面：

\textbf{Lambert 反射定律的操作化。} 多数光学教材中，反射定律以公式形式呈现，学生能够套用计算但未必理解各物理量的量测含义。本实验将 $I_{\mathrm{ref}}=I_{\mathrm{inc}}\rho\cos\theta$ 嵌入传感器通道响应的光谱积分 $V_c(t)=\alpha_c+\beta_c\cos\theta(t)$，使"反射率"从抽象参数转化为可观测的时序信号。

\textbf{光谱加权概念的直观化。} 通过三高斯皮肤反射率模型与 TCS34725 通道响应曲线的卷积计算，不同通道 $\beta_c/\alpha_c$ 比值差异的物理来源——皮肤对红光反射强、对蓝光弱——通过数据变得可见。

\textbf{散粒噪声效应的可测量化。} 学生通过改变环境光强度（暗光→室内常光→台灯→窗边强光），观察识别准确率随 $I_{\mathrm{amb}}$ 增强而下降的趋势，并将实测数据与 $\mathrm{SNR}_c\propto\bar{\rho}\Delta\cos\theta/\sqrt{I_{\mathrm{amb}}}$ 的预测曲线对照，使散粒噪声极限从公式变为可讨论的物理现象。

\subsection{能力培养定位}
\label{sec:adaptability_skills}

实验围绕"建模—仿真—实现—分析"四环节展开，训练三项核心能力：从物理定律出发构建可计算模型的建模能力；将解析仿真与实测波形对照并归因差异的验证能力；将光学、光谱学、信号处理与嵌入式实现贯通使用的整合能力。

\subsection{推广条件与改进方向}
\label{sec:adaptability_promotion}

本实验硬件成本低（TCS34725 模块配合通用 ESP32/Arduino 开发板即可运行），全部器件为市售标准件，软件基于开源平台，不依赖特殊定制设备，具备良好的可复现性。核心参数（皮肤拟合参数、判决阈值）均可根据学习者条件灵活调整，不同肤色的拟合参数差异本身即可作为"通用模型与个体差异"的教学讨论素材。

后续改进可考虑两个方向：一是增加环境光自适应标定环节，使实验在多变光照条件下获得稳定基准数据；二是将实验与生理信号采集相结合，拓展为"手势交互+生理监测"的复合型项目，为不同层次的学习者提供差异化的探索空间。

% =====================================================
\section{结论与局限性}
\label{sec:conclusion}

\subsection{结论}

本文设计并实施了一项基于普通 RGBC 颜色传感器的本科项目式光学实验。该实验把一组原本用于产品研发的 RGBC 手势识别研究，改造为可在 16~32 学时内完成的"建模—仿真—实现—分析"教学项目。学生从 Lambert 反射定律出发，推导四通道探测器读数与入射角之间的关系；用三高斯模型拟合皮肤反射率光谱；实现三阈值判决算法；并将实测准确率衰减与散粒噪声极限 SNR 模型对应。课堂观察与学生反馈表明：本实验能够帮助学生把光度学与光谱响应等抽象概念与一项可触摸的工程应用建立联系。

\subsection{局限性}

本文方法的局限性主要为：（1）RGBC 复用方法在强环境光下散粒噪声主导，SNR 随 $\sqrt{I_{\mathrm{amb}}}$ 退化，准确率受限于式~\eqref{eq:snr} 的物理上界；（2）皮肤反射率个体差异未被显式建模，深色皮肤群体需重新标定 $\delta$ 与 $T_{\mathrm{static}}$；（3）三层阈值依赖启发式选取，跨设备/跨环境泛化能力有限；（4）实测样本量与皮肤异质性覆盖有限，最终投稿版本需扩大样本规模；（5）双击检测的 800 ms 间隔对快速连续操作存在漏检风险。

\subsection{对教师的实施建议}

\begin{enumerate}[label=(\arabic*),leftmargin=2.2em]
  \item \textbf{课前提供 Python 模板与数据集样例。} 多数本科生对光谱积分与曲线拟合的数值实现不熟悉，直接让学生从零开始编程会消耗大量学时。建议教师提供填空式 Python 模板与含 4 类手势、4 种环境光的小样本数据集。
  \item \textbf{课中引导学生对比"规则判决"与"机器学习"。} 在学生完成规则判决算法后，建议教师展示一个简单的机器学习分类器（如 KNN、SVM）在同一数据集上的表现，让学生讨论"为何规则算法在物理含义清晰、阈值稳定时仍优于黑箱模型"。
  \item \textbf{将肤色差异作为光学公平性讨论素材。} 不同肤色被试的准确率差异（约 3--5 个百分点）是本实验最具教学价值的现象之一。教师可借此引导学生从物理量 $\bar{\rho}$ 的角度分析差异来源。
  \item \textbf{允许学生修改阈值并报告 ROC 曲线。} 三阈值（$T_{\mathrm{static}}$、$T_{\mathrm{fast}}$、$\delta$）的调整对准确率的影响是直观的实验素材，建议教师把阈值修改作为可选拓展任务。
\end{enumerate}

\vspace{0.5em}
\noindent\textbf{致谢。} 感谢指导教师的悉心指导与所在单位提供的实验条件支持；感谢 2024 年秋季学期"传感器与检测技术"课程 32 名本科生的实验参与与反馈。

\bibliographystyle{unsrtnat}
\bibliography{references}

\end{document}
